\section{BERT：Encoder-Only 与 Bidirectional Representation}

BERT 展示另一条路线：只使用 encoder block，让每个 token 双向读取上下文，并通过 masked language modeling 学习 representation。它不以开放式自回归生成为主要接口，而是为分类、span prediction、检索和 token labeling 等下游任务提供可 fine-tune 的 contextual embedding。

\subsection{Masked language modeling}

\term{masked language modeling}（MLM，掩码语言建模）随机选择部分输入位置并要求模型预测原 token。由于目标位置被遮蔽，encoder 可以使用左右两侧上下文而不直接读取答案。原始 BERT 还使用 next sentence prediction（NSP）；后续模型对 NSP 的必要性有不同设计，但双向 encoder representation 是更持久的核心。

\lecturefigure{slide-17-bert.jpg}{BERT 用 encoder-only 双向表示与 MLM 支撑下游理解任务。}{官方视频 20:30；Devlin et al. 2018。}

\begin{knowledgebox}{读图：BERT 与 GPT 的差异首先是可见性}
BERT 的 token 可以同时读取左、右上下文，因此适合学习完整句子的 contextual representation；GPT 的 causal token 只能读取左侧前缀，因此天然匹配逐步生成。BERT fine-tuning 常加 task head 并更新参数，GPT prompting 可在统一生成接口中描述任务。两者都来自 Transformer，但 mask、objective 和 product interface 不同。
\end{knowledgebox}

\begin{importantbox}{Encoder-only 与 decoder-only 对照}
\begin{tabular}{L{0.18\textwidth} L{0.36\textwidth} L{0.36\textwidth}}
\toprule
维度 & Encoder-only（BERT） & Decoder-only（GPT） \\
\midrule
可见性 & 双向 self-attention & causal left-to-right attention \\
预训练目标 & masked token prediction & next-token prediction \\
典型输出 & contextual representation & token distribution / generation \\
任务接口 & fine-tuning、linear head、embedding & prompting、completion、fine-tuning \\
优势 & 理解与表示密集任务 & 统一生成与交互接口 \\
\bottomrule
\end{tabular}
\end{importantbox}

\teachervoice{讲者用 BERT 收束导论，强调 encoder block、masked prediction 与下游 fine-tuning。这个顺序很重要：先从原始 encoder-decoder 拆出两个方向，再用 GPT/BERT 说明架构家族，而不是从品牌名称反推机制。}

\subsection{本章小结}

BERT 的核心是双向 encoder representation 和 masked prediction。它与 GPT 共享 self-attention block，却通过不同 mask 和 objective 服务不同任务接口。

